\subsection{Phân biệt hosted service, Hangfire và Quartz.NET}

\textbf{Hangfire} là thư viện .NET dùng để lưu, xếp hàng và thực thi Background Job.
Thông tin job được lưu trong database, nhờ đó job vẫn còn sau khi ứng dụng khởi động
lại. Hangfire hỗ trợ tự động thử lại và cung cấp giao diện theo dõi trạng thái
\cite{hangfire-docs,hangfire-retry}. Hangfire cũng hỗ trợ đăng ký job chạy lặp theo
lịch \cite{hangfire-recurring}.

\textbf{Quartz.NET} là thư viện lập lịch cho .NET. Quartz.NET tách việc cần làm thành
\textbf{job} và thời điểm chạy thành
\textbf{trigger}; công cụ này phù hợp khi bài toán trọng tâm là biểu thức cron
(chuỗi mô tả lịch lặp) hoặc lịch ngày phức tạp
\cite{quartz-jobs-triggers,quartz-job-stores}.

{\scriptsize
\renewcommand{\arraystretch}{1.0}
\begin{longtable}{|>{\raggedright\arraybackslash}p{0.18\textwidth}|
                  >{\raggedright\arraybackslash}p{0.25\textwidth}|
                  >{\raggedright\arraybackslash}p{0.23\textwidth}|
                  >{\raggedright\arraybackslash}p{0.24\textwidth}|}
\caption{So sánh hosted service, Hangfire và Quartz.NET}\label{tab:job-tools}\\
\hline
\textbf{Tiêu chí} & \textbf{Hosted service} & \textbf{Hangfire} & \textbf{Quartz.NET} \\
\hline
Lưu job & Không tự lưu qua lần khởi động lại; ứng dụng phải thiết kế thêm. & Lưu job vào database đã cấu hình. & Có thể lưu trong RAM hoặc database. \\
\hline
Lập lịch & Tự viết vòng lặp, bộ hẹn giờ hoặc hàng đợi. & Hỗ trợ job chạy ngay, chạy trễ và chạy lặp. & Mạnh về trigger, biểu thức cron và lịch ngày. \\
\hline
Nên chọn khi & Background Job nhỏ, cùng tiến trình và chấp nhận tự quản lý độ bền. & Cần lưu job, thử lại và vận hành thuận tiện. & Lịch hoặc quan hệ job--trigger phức tạp. \\
\hline
\end{longtable}
}

Background Job nhỏ và không cần lưu từng job thường chỉ cần hosted service. Khi job không
được phép mất hoặc cần lịch và theo dõi vận hành, mới cần công cụ bên ngoài.
